% !TEX root = ../algebra_02.tex

\section{Линейные отображения.}

\subsection{Внешняя прямая сумма}

\begin{Definition}{Внешняя прямая сумма}{}
	Пусть $V_1, \ldots, V_k$~--- линейные пространства над полем $K$.
	
	\textbf{Внешняя прямая сумма} --- это декартово произведение $V_1 \times \ldots \times V_k$
	с операциями:
	\begin{itemize}
		\item $(v_1, \ldots, v_k) + (v'_1, \ldots, v'_k) = (v_1 + v'_1, \ldots, v_k + v'_k)$
		\item $\lambda\,(v_1, \ldots, v_k) = (\lambda v_1, \ldots, \lambda v_k)$
	\end{itemize}
\end{Definition}

\begin{Statement}{}{}
	$V_1 \times \ldots \times V_k$~--- линейное пространство.
\end{Statement}

\subsubsection*{Связь с внутренней суммой}

Пусть $W_1, \ldots, W_k \subset V$. Тогда существует отображение
\[
\varphi \colon W_1 \times \ldots \times W_k \longrightarrow V,
\qquad
(w_1, \ldots, w_k) \longmapsto w_1 + \ldots + w_k.
\]

\begin{Proposition}{}{}
	\begin{itemize}
		\item $V = W_1 \oplus \ldots \oplus W_k \;\Longleftrightarrow\; \varphi$~--- биекция.
		\item $\varphi$~--- \textbf{всегда} линейное отображение.
	\end{itemize}
\end{Proposition}

\begin{proof}
	\begin{enumerate}
		\item Следует из определения внутренней прямой суммы.
		\item Для $\varphi$ выполнены:
		\begin{itemize}
			\item $\varphi(w + w') = \varphi(w) + \varphi(w')$,
			\item $\varphi(\lambda w) = \lambda\,\varphi(w)$.
		\end{itemize}
	\end{enumerate}
\end{proof}

\begin{Definition}{Изоморфизм линейных пространств}{}
	Пространства $V$ и $W$ над полем $K$ \textbf{изоморфны}, если существует
	биективное линейное отображение $\varphi \colon V \longrightarrow W$.
\end{Definition}

\begin{remark}
	Пока что линейное отображение понимаем как выполнение двух условий:
	\begin{itemize}
		\item $\varphi(x_1 + x_2) = \varphi(x_1) + \varphi(x_2)$,
		\item $\varphi(\lambda x) = \lambda\,\varphi(x)$.
	\end{itemize}
	Точное определение будет введено ниже.
\end{remark}

\begin{Corollary}{}{}
	$V$ изоморфно своей внутренней прямой сумме $W_1 \oplus \ldots \oplus W_k$.
\end{Corollary}

\subsection{Линейные отображения}

\begin{Definition}{Линейное отображение}{}
	Пусть $V$, $W$~--- линейные пространства над полем $K$.
	Отображение $\mathcal{A} \colon V \longrightarrow W$ называется \textbf{линейным}, если
	\[
	\forall\, v_1, v_2 \in V,\quad \forall\, \lambda_1, \lambda_2 \in K \colon
	\qquad
	\mathcal{A}(\lambda_1 v_1 + \lambda_2 v_2) = \lambda_1\,\mathcal{A}v_1 + \lambda_2\,\mathcal{A}v_2.
	\]
\end{Definition}

\begin{Proposition}{}{}
	$\mathcal{A}$~--- линейно $\;\Longleftrightarrow\;$ выполнены оба условия:
	\begin{itemize}
		\item $\mathcal{A}(v_1 + v_2) = \mathcal{A}v_1 + \mathcal{A}v_2$ \quad (\textbf{аддитивность}),
		\item $\mathcal{A}(\lambda v) = \lambda\,\mathcal{A}v$ \quad (\textbf{однородность}).
	\end{itemize}
\end{Proposition}

\begin{proof}
	$(\Rightarrow)$:
	\begin{enumerate}
		\item Аддитивность: положим $\lambda_1 = \lambda_2 = 1$.
		\item Однородность: $\mathcal{A}(\lambda v) = \mathcal{A}(\lambda \cdot v + 0 \cdot 0) = \lambda\,\mathcal{A}v + 0\cdot\mathcal{A}(0) = \lambda\,\mathcal{A}v$.
	\end{enumerate}
	
	$(\Leftarrow)$:
	\[
	\mathcal{A}(\lambda_1 v_1 + \lambda_2 v_2)
	\overset{\text{адд.}}{=} \mathcal{A}(\lambda_1 v_1) + \mathcal{A}(\lambda_2 v_2)
	\overset{\text{одн.}}{=} \lambda_1\,\mathcal{A}v_1 + \lambda_2\,\mathcal{A}v_2.
	\]
\end{proof}

\begin{remark}
	Для любого линейного отображения $\mathcal{A} \colon V \to W$ выполнено $\mathcal{A}(0_V) = 0_W$.
\end{remark}

\begin{remark}
	Существование биективного линейного отображения между $V$ и $W$ равносильно их изоморфизму.
\end{remark}

\begin{Definition}{Множество линейных отображений}{}
	\[
	\mathrm{Hom}(V, W) := \{\,\mathcal{A} \colon V \to W \mid \mathcal{A}\ \text{--- линейное}\,\}.
	\]
\end{Definition}

\begin{Example}{Линейные отображения}{}
	\begin{enumerate}\setlength{\itemsep}{6pt}
		\item Нулевое отображение: $0 \colon V \to W$, $v \mapsto 0$.
		\item Растяжение: $\lambda\,\mathrm{id}_V \colon V \to V$, $v \mapsto \lambda v$, для фиксированного $\lambda \in K$.
		\item Умножение на матрицу. Пусть $A \in K^{m \times n}$:
		\begin{itemize}
			\item $A\cdot \colon K^{n \times p} \to K^{m \times p}$, $B \mapsto AB$,
			\item $\cdot A \colon K^{p \times m} \to K^{p \times n}$, $B \mapsto BA$.
		\end{itemize}
		\item Транспонирование матриц.
		\item Вычисление многочленов:
		\begin{itemize}
			\item Фиксируем $\alpha \in K$: тогда $f \mapsto f(\alpha)$~--- линейное отображение из $K[x]$ в $K$.
			\item Фиксируем $g \in K[x]$: тогда $f \mapsto f \cdot g$~--- линейное отображение из $K[x]$ в $K[x]$.
		\end{itemize}
		\item Дифференцирование: $D \colon f \mapsto f'$. Из линейности дифференцирования
		$(\lambda_1 f_1 + \lambda_2 f_2)' = \lambda_1 f_1' + \lambda_2 f_2'$ следует, что $D \in \mathrm{Hom}(K[x], K[x])$.
	\end{enumerate}
\end{Example}

\begin{Proposition}{Разложение по базису}{}
	Пусть $H = (v_1, \ldots, v_n)$~--- набор векторов в $V$. Определим отображение
	\[
	\varepsilon_H \colon K^n \longrightarrow V,
	\qquad
	\begin{pmatrix} \alpha_1 \\ \vdots \\ \alpha_n \end{pmatrix}
	\longmapsto \alpha_1 v_1 + \ldots + \alpha_n v_n.
	\]
	Тогда:
	\begin{enumerate}\setlength{\itemsep}{7pt}
		\item $\varepsilon_H \in \mathrm{Hom}(K^n, V)$.
		\item $\varepsilon_H$~--- изоморфизм $\;\Leftrightarrow\; H$~--- базис $V$.
	\end{enumerate}
\end{Proposition}

\subsubsection*{Изоморфизм}

Является ли изоморфизм симметричным отношением?

\begin{Proposition}{Симметричность изоморфизма}{}
	\[
	\mathcal{A} \in \mathrm{Hom}(V, W) \text{ --- изоморфизм}
	\;\Longrightarrow\;
	\mathcal{A}^{-1} \in \mathrm{Hom}(W, V).
	\]
\end{Proposition}

\begin{proof}
	Нужно проверить линейность $\mathcal{A}^{-1}$:
	\[
	\mathcal{A}^{-1}(\lambda_1 w_1 + \lambda_2 w_2) \overset{?}{=} \lambda_1\,\mathcal{A}^{-1}w_1 + \lambda_2\,\mathcal{A}^{-1}w_2.
	\]
	Обозначим левую часть $L$, правую~--- $R$. Применим $\mathcal{A}$ к обеим:
	\[
	\mathcal{A}(L) = \lambda_1 w_1 + \lambda_2 w_2,
	\qquad
	\mathcal{A}(R) = \lambda_1\,\mathcal{A}\mathcal{A}^{-1}w_1 + \lambda_2\,\mathcal{A}\mathcal{A}^{-1}w_2 = \lambda_1 w_1 + \lambda_2 w_2.
	\]
	Поскольку $\mathcal{A}$~--- биекция, из $\mathcal{A}(L) = \mathcal{A}(R)$ следует $L = R$.
\end{proof}

\begin{remark}
	Изоморфизм на конечном наборе линейных пространств~--- отношение \textbf{эквивалентности}.
\end{remark}

\begin{Corollary}{}{}
	Для любого пространства $V$ над $K$:
	\[
	\dim V = n \;\Rightarrow\; V \cong K^n.
	\]
\end{Corollary}

\begin{proof}
	Пусть $H$~--- базис $V$. Тогда $\varepsilon_H \colon K^n \to V$~--- изоморфизм, а обратное к нему
	\[
	\varepsilon_H^{-1} \colon V \longrightarrow K^n, \quad v \longmapsto [v]_H
	\]
	сопоставляет каждому вектору его координатный столбец в базисе $H$.
\end{proof}

\subsubsection*{Универсальное свойство базиса}

\begin{Proposition}{}{}
	Пусть $V$, $W$~--- пространства над $K$, $\;(e_1, \ldots, e_n)$~--- базис $V$,
	и $w_1, \ldots, w_n \in W$~--- произвольный набор векторов. Тогда
	\[
	\exists!\; \mathcal{A} \in \mathrm{Hom}(V, W) \colon \mathcal{A}(e_j) = w_j \quad \forall\, j.
	\]
\end{Proposition}

\begin{proof}
	\begin{enumerate}
		\item \textit{Существование.} Определим $\mathcal{A}$ формулой
		\[
		\mathcal{A}(v) = \varepsilon_H\!\left([v]_E\right),
		\]
		где $E = (e_1, \ldots, e_n)$~--- базис $V$, $H = (w_1, \ldots, w_n)$. Явно:
		\[
		\mathcal{A}(\alpha_1 e_1 + \ldots + \alpha_n e_n) = \alpha_1 w_1 + \ldots + \alpha_n w_n.
		\]
		В частности, $\mathcal{A}(e_j) = w_j$. Линейность $\mathcal{A}$ следует из линейности $\varepsilon_H$.
		
		\item \textit{Единственность.} Если $\mathcal{A}' \in \mathrm{Hom}(V, W)$ также удовлетворяет $\mathcal{A}'(e_j) = w_j$, то для любого $v = \alpha_1 e_1 + \ldots + \alpha_n e_n$:
		\[
		\mathcal{A}'v = \alpha_1\,\mathcal{A}'e_1 + \ldots + \alpha_n\,\mathcal{A}'e_n = \alpha_1 w_1 + \ldots + \alpha_n w_n = \mathcal{A}v.
		\]
		Следовательно, $\mathcal{A}' = \mathcal{A}$.
	\end{enumerate}
\end{proof}

\subsection{Ядро и образ линейного отображения}

Пусть $\mathcal{A} \in \mathrm{Hom}(V, W)$.

\begin{Definition}{Ядро линейного отображения}{}
	\[
	\Ker\mathcal{A} := \{\, v \in V \mid \mathcal{A}v = 0 \,\}.
	\]
\end{Definition}

\begin{Definition}{Образ линейного отображения}{}
	\[
	\Im\mathcal{A} := \{\, w \in W \mid \exists\, v \in V \colon \mathcal{A}v = w \,\}.
	\]
\end{Definition}

\begin{Proposition}{}{}
	\begin{enumerate}
		\item $\Ker\mathcal{A}$~--- линейное подпространство в $V$.
		\item $\Im\mathcal{A}$~--- линейное подпространство в $W$.
	\end{enumerate}
\end{Proposition}

\begin{proof}
	Напомним, что подпространство $V' \subset V$ определяется двумя условиями:
	$0 \in V'$ и $\lambda_1 v_1 + \lambda_2 v_2 \in V'$ для любых $v_1, v_2 \in V'$, $\lambda_1, \lambda_2 \in K$.
	
	\begin{enumerate}
		\item \textit{Ядро.}
		\begin{itemize}
			\item $\mathcal{A}(0) = 0$, поэтому $0 \in \Ker\mathcal{A}$.
			\item Пусть $v_1, v_2 \in \Ker\mathcal{A}$. Тогда
			\[
			\mathcal{A}(\lambda_1 v_1 + \lambda_2 v_2)
			= \lambda_1\,\mathcal{A}v_1 + \lambda_2\,\mathcal{A}v_2
			= \lambda_1 \cdot 0 + \lambda_2 \cdot 0 = 0,
			\]
			то есть $\lambda_1 v_1 + \lambda_2 v_2 \in \Ker\mathcal{A}$.
		\end{itemize}
		\item \textit{Образ.}
		\begin{itemize}
			\item $\mathcal{A}(0) = 0$, поэтому $0 \in \Im\mathcal{A}$.
			\item Пусть $w_1, w_2 \in \Im\mathcal{A}$, то есть $\mathcal{A}v_1 = w_1$, $\mathcal{A}v_2 = w_2$ для некоторых $v_1, v_2 \in V$. Тогда
			\[
			\mathcal{A}(\lambda_1 v_1 + \lambda_2 v_2)
			= \lambda_1\,\mathcal{A}v_1 + \lambda_2\,\mathcal{A}v_2
			= \lambda_1 w_1 + \lambda_2 w_2,
			\]
			откуда $\lambda_1 w_1 + \lambda_2 w_2 \in \Im\mathcal{A}$.
		\end{itemize}
	\end{enumerate}
\end{proof}

\begin{Proposition}{Теорема о ранге и дефекте}{}
	Пусть $\dim V = n < \infty$. Тогда $\Ker\mathcal{A}$ и $\Im\mathcal{A}$~--- конечномерные, причём
	\[
	\dim\Ker\mathcal{A} + \dim\Im\mathcal{A} = n.
	\]
	Величина $\dim\Ker\mathcal{A} =: \delta\mathcal{A}$ называется \textbf{дефектом} $\mathcal{A}$,
	а $\dim\Im\mathcal{A} =: \mathrm{rk}\,\mathcal{A}$~--- \textbf{рангом} $\mathcal{A}$.
\end{Proposition}

\begin{proof}
	Из $\Ker\mathcal{A} \subset V$ следует $d := \dim\Ker\mathcal{A} < \infty$.
	Пусть $e_1, \ldots, e_d$~--- базис ядра. Дополним его до базиса $V$:
	\[
	e_1, \ldots, e_d,\; e_{d+1}, \ldots, e_n.
	\]
	Утверждаем, что $\mathcal{A}e_{d+1}, \ldots, \mathcal{A}e_n$~--- базис $\Im\mathcal{A}$.
	
	\textit{Система образующих.} Для любого $w \in \Im\mathcal{A}$ найдём $v \in V$ с $\mathcal{A}v = w$:
	\[
	w = \mathcal{A}(\alpha_1 e_1 + \ldots + \alpha_n e_n)
	= \underbrace{\mathcal{A}(\alpha_1 e_1 + \ldots + \alpha_d e_d)}_{\in\,\Ker\mathcal{A},\;\text{т.е. }=\,0}
	+ \alpha_{d+1}\,\mathcal{A}e_{d+1} + \ldots + \alpha_n\,\mathcal{A}e_n.
	\]
	Следовательно, $w = \alpha_{d+1}\,\mathcal{A}e_{d+1} + \ldots + \alpha_n\,\mathcal{A}e_n$.
	
	\textit{Линейная независимость.} Пусть
	\[
	(\alpha_{d+1} - \alpha_{d+1}')\,\mathcal{A}e_{d+1} + \ldots + (\alpha_n - \alpha_n')\,\mathcal{A}e_n = 0.
	\]
	Тогда $\mathcal{A}\bigl((\alpha_{d+1}-\alpha_{d+1}')e_{d+1} + \ldots + (\alpha_n-\alpha_n')e_n\bigr) = 0$,
	то есть этот вектор лежит в $\Ker\mathcal{A} = \mathrm{Lin}(e_1, \ldots, e_d)$.
	Но тогда $e_{d+1}, \ldots, e_n$ линейно выражались бы через $e_1, \ldots, e_d$~--- противоречие с тем,
	что всё вместе они образуют базис $V$. Значит, все коэффициенты равны нулю.
	
	Итого, $\dim\Im\mathcal{A} = n - d$.
\end{proof}

\subsubsection*{ОСЛУ и ранг отображения}

\begin{Theorem}{Множество решений ОСЛУ}{}
	Пусть $A \in K^{m \times n}$. Тогда множество решений однородной системы линейных уравнений
	$Ax = 0$ является линейным подпространством $K^n$ размерности $n - \mathrm{rk}\,A$.
\end{Theorem}

\begin{proof}
	Рассмотрим линейное отображение $\mathcal{A} = (A\,\cdot) \colon K^n \to K^m$, $x \mapsto Ax$.
	Тогда множество решений $H = \Ker\mathcal{A}$.
	
	Найдём $\dim\Im\mathcal{A}$. Из формулы матричного умножения:
	\[
	Ax = x_1\,A_{*1} + \ldots + x_n\,A_{*n},
	\]
	где $A_{*j}$~--- $j$-й столбец матрицы $A$. Следовательно,
	$\Im\mathcal{A} = \mathrm{Lin}(A_{*1}, \ldots, A_{*n})$,
	и по определению $\dim\Im\mathcal{A} = \mathrm{rk}\,A$.
	
	По теореме о ранге и дефекте:
	\[
	\dim H = \dim\Ker\mathcal{A} = n - \dim\Im\mathcal{A} = n - \mathrm{rk}\,A.
	\]
\end{proof}

\begin{Corollary}{}{}
	Однородная СЛУ определена (имеет единственное решение) тогда и только тогда, когда $\mathrm{rk}\,A = n$.
\end{Corollary}

\begin{proof}
	\[
	\exists!\text{ решение} \;\Leftrightarrow\; \dim H = 0 \;\Leftrightarrow\; n - \mathrm{rk}\,A = 0
	\;\Leftrightarrow\; \mathrm{rk}\,A = n.
	\]
\end{proof}

\subsubsection*{Инъективность и сюръективность линейного отображения}

\textbf{Образ} отображения~--- мера его \textbf{сюръективности};
\textbf{ядро}~--- мера его \textbf{инъективности}.

\begin{Proposition}{}{}
	$\mathcal{A} \in \mathrm{Hom}(V, W)$ инъективно $\;\Longleftrightarrow\; \Ker\mathcal{A} = \{0\}$.
\end{Proposition}

\begin{proof}
	\begin{itemize}
		\item $(\Rightarrow)$: При инъективности $|\mathcal{A}^{-1}(0)| \le 1$.
		Так как $\mathcal{A}(0) = 0$, получаем $\Ker\mathcal{A} = \{0\}$.
		\item $(\Leftarrow)$: Пусть $\mathcal{A}v_1 = \mathcal{A}v_2$. Тогда
		$\mathcal{A}(v_1 - v_2) = 0$, то есть $v_1 - v_2 \in \Ker\mathcal{A} = \{0\}$,
		откуда $v_1 = v_2$.
	\end{itemize}
\end{proof}

\begin{Definition}{}{}
	\begin{itemize}
		\item \textbf{Мономорфизм}~--- инъективное линейное отображение.
		\item \textbf{Эпиморфизм}~--- сюръективное линейное отображение.
		\item \textbf{Изоморфизм}~--- биективное линейное отображение.
	\end{itemize}
\end{Definition}

\begin{Proposition}{}{}
	Пусть $\mathcal{A} \in \mathrm{Hom}(V, W)$, и $(e_1, \ldots, e_n)$~--- базис $V$. Тогда:
	\begin{enumerate}
		\item $\mathcal{A}$~--- инъективно $\;\Leftrightarrow\; \mathcal{A}e_1, \ldots, \mathcal{A}e_n$~--- \textbf{линейно независимы} в $W$.
		\item $\mathcal{A}$~--- сюръективно $\;\Leftrightarrow\; \mathcal{A}e_1, \ldots, \mathcal{A}e_n$~--- \textbf{система образующих} $W$.
		\item $\mathcal{A}$~--- биективно $\;\Leftrightarrow\; \mathcal{A}e_1, \ldots, \mathcal{A}e_n$~--- \textbf{базис} $W$.
	\end{enumerate}
\end{Proposition}

\begin{proof}
	\begin{enumerate}
		\item
		\begin{itemize}
			\item $(\Rightarrow)$: Пусть $\lambda_1\,\mathcal{A}e_1 + \ldots + \lambda_n\,\mathcal{A}e_n = 0$.
			Тогда $\mathcal{A}(\lambda_1 e_1 + \ldots + \lambda_n e_n) = 0$,
			и из $\Ker\mathcal{A} = \{0\}$ следует $\lambda_1 e_1 + \ldots + \lambda_n e_n = 0$,
			откуда $\lambda_1 = \ldots = \lambda_n = 0$.
			\item $(\Leftarrow)$: Пусть $\mathcal{A}v = 0$, $v = \lambda_1 e_1 + \ldots + \lambda_n e_n$.
			Тогда $\lambda_1\,\mathcal{A}e_1 + \ldots + \lambda_n\,\mathcal{A}e_n = 0$,
			и из линейной независимости $\lambda_1 = \ldots = \lambda_n = 0$, то есть $v = 0$.
		\end{itemize}
		\item
		\begin{itemize}
			\item $(\Rightarrow)$: Для любого $w \in W$ найдём $v \in V$ с $\mathcal{A}v = w$.
			Разложив $v$ по базису:
			\[
			w = \mathcal{A}(\lambda_1 e_1 + \ldots + \lambda_n e_n) = \lambda_1\,\mathcal{A}e_1 + \ldots + \lambda_n\,\mathcal{A}e_n.
			\]
			\item $(\Leftarrow)$: Для любого $w \in W$ запишем $w = \lambda_1\,\mathcal{A}e_1 + \ldots + \lambda_n\,\mathcal{A}e_n = \mathcal{A}(\lambda_1 e_1 + \ldots + \lambda_n e_n) \in \Im\mathcal{A}$.
		\end{itemize}
		\item Следует из пунктов 1 и 2.
	\end{enumerate}
\end{proof}

\begin{Statement}{}{}
	Пусть $\dim V = \dim W < \infty$ и $\mathcal{A} \in \mathrm{Hom}(V, W)$.
	Тогда следующие условия равносильны:
	\begin{enumerate}
		\item $\mathcal{A}$~--- инъективно,
		\item $\mathcal{A}$~--- сюръективно,
		\item $\mathcal{A}$~--- биективно.
	\end{enumerate}
\end{Statement}

\noindent(Аналог принципа Дирихле для линейных отображений.)

\begin{proof}
	По теореме о ранге и дефекте:
	\[
	\dim\Ker\mathcal{A} + \dim\Im\mathcal{A} = n.
	\]
	Поэтому
	\[
	\Ker\mathcal{A} = \{0\}
	\;\Leftrightarrow\; \dim\Im\mathcal{A} = n
	\;\Leftrightarrow\; \Im\mathcal{A} = W. \qedhere
	\]
\end{proof}

\subsubsection*{Матрица линейного отображения}

Пусть $\dim V = n$, $\dim W = m$, $\mathcal{A} \in \mathrm{Hom}(V, W)$,
$E = (e_1, \ldots, e_n)$~--- базис $V$, $F = (f_1, \ldots, f_m)$~--- базис $W$.

\begin{Definition}{Матрица линейного отображения}{}
	\textbf{Матрицей отображения} $\mathcal{A}$ в базисах $E$ и $F$ называется матрица
	\[
	[\mathcal{A}]_{E,F} = \bigl([\mathcal{A}e_1]_F \;\Big|\; \ldots \;\Big|\; [\mathcal{A}e_n]_F\bigr) \in K^{m \times n},
	\]
	столбцами которой служат координатные столбцы образов базисных векторов $e_j$ в базисе $F$.
	
	Условно записывают:
	\[
	\mathcal{A}(E) = F \cdot [\mathcal{A}]_{E,F},
	\]
	понимая правую часть как строку векторов, каждый из которых~--- линейная комбинация $f_i$
	с коэффициентами из соответствующего столбца матрицы:
	\[
	F \cdot [\mathcal{A}]_{E,F}
	:=
	\left(
	\sum_{i=1}^m a_{i1} f_i,\;
	\ldots,\;
	\sum_{i=1}^m a_{in} f_i
	\right).
	\]
\end{Definition}

\begin{remark}
	Столбец $[\mathcal{A}e_j]_F$~--- это координатный столбец вектора $\mathcal{A}e_j \in W$ в базисе $F$.
\end{remark}